% Part: computability
% Chapter: tm-computations
% Section: church-turing-thesis

\documentclass[../../../include/open-logic-section]{subfiles}

\begin{document}

\olfileid{tur}{mac}{ctt}
\olsection{চার্চ--টুরিং থিসিস}

টুরিং যন্ত্রকে কার্যকর পদ্ধতির ধারণার একটি নির্ভুল প্রতিস্থাপক হিসেবে
ধরা হয়। টুরিং মনে করতেন, কার্যকর পদ্ধতি ও টুরিং যন্ত্র—দুটি ধারণাই যিনি
বুঝেছেন, তাঁর স্বজ্ঞা হবে যে কার্যকর পদ্ধতিতে করা যায় এমন সবকিছুই টুরিং
যন্ত্রে করা যায়। কার্যকর পদ্ধতির ধারণার অন্য যত নির্ভুল প্রতিস্থাপক
প্রস্তাব করা হয়েছে, সেগুলি শেষ পর্যন্ত টুরিং যন্ত্রের ধারণার সঙ্গে
বহির্বিস্তারণগতভাবে সমতুল্য—অর্থাৎ তারা ঠিক একই অপেক্ষক-সেট গণনা করতে
পারে—এই তথ্যটি দাবিটির পক্ষে সমর্থন দেয়। দাবিটিকে \emph{চার্চ--টুরিং
থিসিস} বলা হয়।

\begin{defn}[চার্চ--টুরিং থিসিস]
\emph{চার্চ--টুরিং থিসিস} বলে, কার্যকর পদ্ধতিতে গণনসাধ্য সবকিছুই
টুরিং-গণনসাধ্য।
\end{defn}

চার্চ--টুরিং থিসিসকে দুই ভাবে কাজে লাগানো হয়। প্রথম ব্যবহারটি
পরিশ্রম বাঁচানোর অজুহাত। ধরা যাক, কোনো কিছু গণনা করার কার্যকর পদ্ধতির
একটি বিবরণ, যেমন ``ছদ্ম-কোডে'', আমাদের আছে। তাহলে যন্ত্রটি বাস্তবে নির্মাণ
না করেও একই অপেক্ষক কোনো টুরিং যন্ত্রে গণনা করা যায়—এই দাবির সমর্থনে
চার্চ--টুরিং থিসিস আহ্বান করা যায়।

চার্চ--টুরিং থিসিসের অন্য ব্যবহারটি দার্শনিকভাবে বেশি আকর্ষণীয়।
দেখানো যায়, এমন অপেক্ষক আছে যা টুরিং যন্ত্রে গণনা করা যায় না। এখান থেকে
চার্চ--টুরিং থিসিস ব্যবহার করে সিদ্ধান্ত নেওয়া যায় যে কোনো পদ্ধতিতেই
সেটিকে কার্যকরভাবে গণনা করা যায় না। কারণ, এমন কোনো পদ্ধতি থাকলে
চার্চ--টুরিং থিসিস অনুযায়ী তার জন্য একটি টুরিং যন্ত্রও থাকত। অতএব
কোনো অপেক্ষক গণনা করে এমন টুরিং যন্ত্র নেই বলে প্রমাণ করতে পারলে তার কোনো
কার্যকর পদ্ধতিও থাকতে পারে না। বিশেষত, তথাকথিত থামা সমস্যা শুধু টুরিং
যন্ত্রে সমাধান করা যায় না তা-ই নয়, আদৌ কার্যকরভাবে সমাধান করা যায়
না—এই দাবি করতেও চার্চ--টুরিং থিসিস আহ্বান করা হয়।

\end{document}
